\section{Results}

\subsection{Data Description}

The analysis is based on a synthetic student-level dataset containing 60 observations. The dataset was constructed for the purpose of this study and contains information on study time, attendance, previous academic performance, access to learning resources, and current academic performance.

Academic performance is treated as the dependent variable, while study time, attendance, previous performance, and access to learning resources are treated as explanatory variables. The dataset therefore provides the variables required to investigate the statistical relationships specified in the methodology.

\subsection{Descriptive Statistics}

Descriptive statistics were calculated to summarise the distribution of the variables included in the analysis. Table~\ref{tab:descriptive} presents the mean, standard deviation, minimum, and maximum values.

\begin{table}[htbp]
\centering
\caption{Descriptive statistics of study variables}
\label{tab:descriptive}
\begin{tabular}{lrrrr}
\toprule
Variable & Mean & Standard Deviation & Minimum & Maximum \\
\midrule
Study Time & 9.15 & 4.23 & 2 & 15 \\
Attendance & 84.85 & 11.17 & 65 & 100 \\
Previous Performance & 70.33 & 12.69 & 46 & 90 \\
Learning Resources & 3.35 & 1.54 & 1 & 5 \\
Academic Performance & 57.03 & 8.42 & 40.0 & 76.5 \\
\bottomrule
\end{tabular}
\end{table}

\subsection{Exploratory Data Analysis}

Exploratory data analysis was conducted to examine the relationships between the explanatory variables and academic performance. The analysis provides an initial assessment of whether differences in study time, attendance, previous academic performance, and access to learning resources are associated with differences in current academic performance.

The patterns observed during this stage provide an initial basis for the subsequent correlation analysis and multiple linear regression model. In particular, the analysis considers whether students with higher levels of study time, attendance, previous performance, and access to learning resources tend to record higher academic performance.

\subsection{Correlation Analysis}

The Pearson correlation coefficient is used to measure the strength and direction of the linear relationship between academic performance and each explanatory variable. The correlation coefficient is calculated as

\begin{equation}
r_{XY} =
\frac{\sum_{i=1}^{n}(X_i-\bar{X})(Y_i-\bar{Y})}
{\sqrt{\sum_{i=1}^{n}(X_i-\bar{X})^2
\sum_{i=1}^{n}(Y_i-\bar{Y})^2}}.
\end{equation}

For the 60 student observations, the mean academic performance is

\begin{equation}
\bar{Y} = \frac{\sum Y_i}{60} = 57.0267.
\end{equation}

The sum of squared deviations of academic performance is

\begin{equation}
\sum_{i=1}^{60}(Y_i-\bar{Y})^2 = 4185.1173.
\end{equation}

\subsubsection{Study Time and Academic Performance}

For Study Time, the mean is $\bar{X}=9.15$. The required components of the correlation calculation are

\begin{align*}
\sum(X_i-\bar{X})(Y_i-\bar{Y}) &= 1314.36,\\
\sum(X_i-\bar{X})^2 &= 1057.65.
\end{align*}

Therefore,

\begin{equation}
r =
\frac{1314.36}
{\sqrt{(1057.65)(4185.1173)}}
=0.6247\approx\mathbf{0.625}.
\end{equation}

\subsubsection{Attendance and Academic Performance}

For Attendance, the mean is $\bar{X}=84.85$. The required components are

\begin{align*}
\sum(X_i-\bar{X})(Y_i-\bar{Y}) &= 831.44,\\
\sum(X_i-\bar{X})^2 &= 7363.65.
\end{align*}

Therefore,

\begin{equation}
r =
\frac{831.44}
{\sqrt{(7363.65)(4185.1173)}}
=0.1498\approx\mathbf{0.150}.
\end{equation}

\subsubsection{Previous Performance and Academic Performance}

For Previous Performance, the mean is $\bar{X}=70.3333$. The required components are

\begin{align*}
\sum(X_i-\bar{X})(Y_i-\bar{Y}) &= 3212.3667,\\
\sum(X_i-\bar{X})^2 &= 9495.3333.
\end{align*}

Therefore,

\begin{equation}
r =
\frac{3212.3667}
{\sqrt{(9495.3333)(4185.1173)}}
=0.5096\approx\mathbf{0.510}.
\end{equation}

\subsubsection{Learning Resources and Academic Performance}

For Learning Resources, the mean is $\bar{X}=3.35$. The required components are

\begin{align*}
\sum(X_i-\bar{X})(Y_i-\bar{Y}) &= 83.44,\\
\sum(X_i-\bar{X})^2 &= 139.65.
\end{align*}

Therefore,

\begin{equation}
r =
\frac{83.44}
{\sqrt{(139.65)(4185.1173)}}
=0.1091\approx\mathbf{0.109}.
\end{equation}

The resulting correlation coefficients are summarised in Table~\ref{tab:correlation}.

\begin{table}[htbp]
\centering
\caption{Correlation between explanatory variables and academic performance}
\label{tab:correlation}
\begin{tabular}{lr}
\toprule
Variable & Correlation with Academic Performance ($r$) \\
\midrule
Study Time & 0.625 \\
Attendance & 0.150 \\
Previous Performance & 0.510 \\
Learning Resources & 0.109 \\
\bottomrule
\end{tabular}
\end{table}


\subsection{Regression Model Calculation}

The multiple linear regression model specified in Equation~(1) is applied to the 60 student observations. For clarity, Equation~(1) is recalled below:

\begin{equation}
Y_i = \beta_0 + \beta_1 X_{1i} + \beta_2 X_{2i}
+ \cdots + \beta_p X_{pi} + \varepsilon_i.
\tag{1}
\end{equation}

For this study, the explanatory variables are defined as

\begin{align*}
X_1 &= \text{Study Time},\\
X_2 &= \text{Attendance},\\
X_3 &= \text{Previous Performance},\\
X_4 &= \text{Learning Resources}.
\end{align*}

Therefore, Equation~(1) becomes

\begin{equation}
Y_i = \beta_0 + \beta_1 X_{1i} + \beta_2 X_{2i}
+ \beta_3 X_{3i} + \beta_4 X_{4i} + \varepsilon_i.
\end{equation}

The regression coefficients are estimated using the ordinary least squares method:

\begin{equation}
\hat{\boldsymbol{\beta}}
=
(X^{T}X)^{-1}X^{T}Y.
\end{equation}

For the 60 observations in the dataset, the design matrix $X$ contains a column of ones for the intercept and four columns corresponding to Study Time, Attendance, Previous Performance, and Learning Resources. The required matrix products are

\begin{equation}
X^{T}X =
\begin{bmatrix}
60 & 549 & 5091 & 4220 & 201\\
549 & 6081 & 46503 & 38277 & 1758\\
5091 & 46503 & 439335 & 358413 & 16856\\
4220 & 38277 & 358413 & 306302 & 14349\\
201 & 1758 & 16856 & 14349 & 813
\end{bmatrix}
\end{equation}

and

\begin{equation}
X^{T}Y =
\begin{bmatrix}
3421.6\\
32622.0\\
291154.2\\
243864.9\\
11545.8
\end{bmatrix}.
\end{equation}

Therefore,

\begin{equation}
\hat{\boldsymbol{\beta}}
=
(X^{T}X)^{-1}X^{T}Y
=
\begin{bmatrix}
2.759\\
1.452\\
0.141\\
0.360\\
1.096
\end{bmatrix}.
\end{equation}

Thus, the fitted regression equation is

\begin{equation}
\widehat{Y}
=
2.759
+1.452X_1
+0.141X_2
+0.360X_3
+1.096X_4.
\end{equation}

The coefficient of determination, $R^2$, is calculated as

\begin{equation}
R^2
=
1-\frac{SS_{\mathrm{res}}}{SS_{\mathrm{tot}}},
\end{equation}

where $SS_{\mathrm{res}}$ is the residual sum of squares and
$SS_{\mathrm{tot}}$ is the total sum of squares.

For the fitted model,

\begin{align*}
SS_{\mathrm{tot}} &= 4185.1173,\\
SS_{\mathrm{res}} &= 911.3299.
\end{align*}
Therefore,

\begin{equation}
R^2
=
1-\frac{911.3299}{4185.1173}
=
0.7822
\approx\mathbf{0.782}.
\end{equation}

The adjusted coefficient of determination accounts for the number of explanatory variables in the model. It is calculated as

\begin{equation}
\bar{R}^{2}
=
1-
\left(1-R^2\right)
\frac{n-1}{n-p-1},
\end{equation}

where $n=60$ is the number of observations and $p=4$ is the number of explanatory variables.

Thus,

\begin{equation}
\bar{R}^{2}
=
1-(1-0.7822)
\frac{59}{55}
=
0.7664
\approx\mathbf{0.766}.
\end{equation}

For example, for a student with 10 hours of study time, 90\% attendance, a previous performance score of 70, and a learning-resource score of 4, the predicted academic performance is

\begin{align*}
\widehat{Y}
&=2.759+(1.452)(10)+(0.141)(90)\\
&\quad +(0.360)(70)+(1.096)(4)\\
&=59.55.
\end{align*}

The fitted coefficients and their statistical significance are presented in Table~\ref{tab:regression}.

\begin{table}[htbp]
\centering
\caption{Multiple linear regression results}
\label{tab:regression}
\begin{tabular}{lrr}
\toprule
Variable & Coefficient & $p$-value \\
\midrule
Intercept & 2.759 & 0.616 \\
Study Time & 1.452 & $<0.001$ \\
Attendance & 0.141 & 0.005 \\
Previous Performance & 0.360 & $<0.001$ \\
Learning Resources & 1.096 & 0.004 \\
\midrule
$R^2$ & 0.782 & \\
Adjusted $R^2$ & 0.766 & \\
F-statistic & 49.394 & $<0.001$ \\
\bottomrule
\end{tabular}
\end{table}


\subsection{Model Diagnostics and Validation}

The fitted model was checked by examining the residuals, which represent the differences between the observed and predicted academic performance values. Each residual is calculated as

\begin{equation}
e_i = Y_i-\widehat{Y}_i.
\end{equation}

Across the 60 observations, the residuals ranged from approximately $-9.663$ to $9.089$, while their mean was approximately zero. The residuals were distributed around zero without a clear systematic pattern.

These checks provide a simple assessment of the fitted model and do not indicate a substantial problem with the model's fit to the synthetic dataset.

\subsection{Interpretation of Regression Results}

The regression results indicate that all four explanatory variables have positive relationships with academic performance within the fitted model.

Study Time has a coefficient of 1.452 and a $p$-value of less than 0.001. Holding the other variables constant, a one-unit increase in study time is associated with an estimated 1.452-unit increase in academic performance.

Attendance has a coefficient of 0.141 and a $p$-value of 0.005. Holding the other variables constant, a one-unit increase in attendance is associated with an estimated 0.141-unit increase in academic performance.

Previous Performance has a coefficient of 0.360 and a $p$-value of less than 0.001. Holding the other variables constant, a one-unit increase in previous performance is associated with an estimated 0.360-unit increase in academic performance.

Learning Resources has a coefficient of 1.096 and a $p$-value of 0.004. Holding the other variables constant, a one-unit increase in the learning-resource score is associated with an estimated 1.096-unit increase in academic performance.

All four explanatory variables have $p$-values below the 5\% significance level. This indicates that each variable provides statistically significant evidence of an association with academic performance within the fitted model.

The model has an $R^2$ value of 0.782, indicating that 78.2\% of the variation in academic performance in the synthetic dataset is explained jointly by the four explanatory variables. The adjusted $R^2$ is 0.766, showing that the model retains substantial explanatory power after accounting for the number of explanatory variables.

Overall, the results suggest that study time, attendance, previous performance, and access to learning resources are useful predictors of academic performance within this synthetic dataset. The fitted model therefore provides a practical statistical framework for examining the factors considered in this study.